\documentclass[aspectratio=169]{beamer}

\usepackage[utf8]{inputenc}
\usepackage[T1]{fontenc}
\usepackage{amsmath,amssymb}
\usepackage{booktabs}

\usetheme{Madrid}
\usecolortheme{default}

\definecolor{BootcampBlue}{HTML}{1A3C68}
\definecolor{BootcampGold}{HTML}{C69214}
\definecolor{BootcampRed}{HTML}{9A1B1E}
\definecolor{BootcampGreen}{HTML}{2F7D32}

\setbeamercolor{structure}{fg=BootcampBlue}
\setbeamercolor{frametitle}{fg=white,bg=BootcampBlue}
\setbeamercolor{title}{fg=white,bg=BootcampBlue}
\setbeamercolor{block title}{fg=white,bg=BootcampBlue}
\setbeamercolor{block body}{fg=black,bg=BootcampBlue!6}
\setbeamercolor{alerted text}{fg=BootcampRed}

\newenvironment{resultblock}[1]{%
  \setbeamercolor{block title}{fg=black,bg=BootcampGold}%
  \setbeamercolor{block body}{fg=black,bg=BootcampGold!10}%
  \begin{block}{#1}}{\end{block}}

\newcommand{\R}{\mathbb{R}}
\newcommand{\M}{\mathcal{M}}
\newcommand{\rank}{\operatorname{rank}}
\newcommand{\im}{\operatorname{Im}}
\newcommand{\tr}{\operatorname{tr}}
\newcommand{\diag}{\operatorname{diag}}
\newcommand{\Cov}{\operatorname{Cov}}
\newcommand{\Var}{\operatorname{Var}}

\title[Session 1 -- Linear Algebra]{Summer Bootcamp\\Session 1 -- Linear Algebra}
\subtitle{Matrices, linear systems, diagonalization, and statistical applications}
\author{Nathan Sauldubois}
\institute{NYU Finance and Risk Engineering}
\date{August 2026}

\AtBeginSection[]{
\begin{frame}{Section Roadmap}
  \small
  \tableofcontents[currentsection,sectionstyle=show/hide,
    subsectionstyle=show/show/hide]
\end{frame}
}

\begin{document}

\begin{frame}
  \titlepage
\end{frame}

\begin{frame}{Session Roadmap}
\begin{block}{Learning goals}
By the end of the session, students should be able to:
\begin{itemize}
  \item manipulate matrices and check dimension compatibility;
  \item write and solve linear systems in matrix form;
  \item interpret image, kernel, rank, invertibility, and transposition;
  \item compute eigenvalues and use diagonalization;
  \item connect symmetric matrices to covariance, portfolio variance, and PCA.
\end{itemize}
\end{block}
\end{frame}

\section{Algebraic Operations}

\begin{frame}[shrink=5]{Matrices}
\begin{block}{Definition}
A matrix is a rectangular array of real numbers:
\[
A=(a_{ij})_{1\le i\le n,\ 1\le j\le m}\in \M_{n,m}(\R).
\]
It can be interpreted as a linear transformation from $\R^m$ to $\R^n$.
\end{block}
\pause
\begin{exampleblock}{Examples}
\[
 A=\begin{pmatrix}1&2&3\\4&5&6\end{pmatrix}\in \M_{2,3}(\R),
 \qquad
 x=\begin{pmatrix}7\\8\\9\end{pmatrix}\in \R^3.
\]
The dimensions match, so
\[
Ax=\begin{pmatrix}50\\122\end{pmatrix}\in\R^2.
\]
\end{exampleblock}
\end{frame}

\begin{frame}{Addition and Scalar Multiplication}
Let $A=(a_{ij})$ and $B=(b_{ij})$ have the same size. Then
\[
A+B=(a_{ij}+b_{ij}),\qquad
\lambda A=(\lambda a_{ij}).
\]
\pause
\begin{exampleblock}{Example}
\[
A=\begin{pmatrix}1&-2&3\\0&4&-1\end{pmatrix},
\quad
B=\begin{pmatrix}2&0&-1\\-3&1&5\end{pmatrix}.
\]
\[
A-2B=
\begin{pmatrix}
-3&-2&5\\
6&2&-11
\end{pmatrix}.
\]
\end{exampleblock}
\end{frame}

\begin{frame}{Matrix Multiplication}
If
\[
A\in \M_{n,m}(\R),\qquad B\in \M_{m,p}(\R),
\]
then $AB\in \M_{n,p}(\R)$ and
\[
(AB)_{ij}=\sum_{k=1}^m a_{ik}b_{kj}.
\]
\pause
\begin{block}{Dimension checks}
\[
(2\times 3)(3\times 1)\to 2\times 1,\qquad
(4\times 2)(2\times 5)\to 4\times 5.
\]
\[
(2\times 3)(2\times 5)\quad\text{is not defined.}
\]
\end{block}
\end{frame}

\begin{frame}{Matrix Multiplication: Worked Example}
\[
A=\begin{pmatrix}
1&2&-1\\
0&3&4
\end{pmatrix},
\qquad
B=\begin{pmatrix}
2&1\\
-1&0\\
3&2
\end{pmatrix}.
\]
\pause
\[
AB=
\begin{pmatrix}
1\cdot2+2(-1)+(-1)3 & 1\cdot1+2\cdot0+(-1)2\\
0\cdot2+3(-1)+4\cdot3 & 0\cdot1+3\cdot0+4\cdot2
\end{pmatrix}
=
\begin{pmatrix}-3&-1\\9&8\end{pmatrix}.
\]
\end{frame}

\begin{frame}{Order Matters}
\[
A=\begin{pmatrix}1&2\\0&1\end{pmatrix},
\qquad
B=\begin{pmatrix}0&1\\1&0\end{pmatrix}.
\]
\pause
\[
AB=\begin{pmatrix}2&1\\1&0\end{pmatrix},
\qquad
BA=\begin{pmatrix}0&1\\1&2\end{pmatrix}.
\]
\pause
\[
\boxed{AB\neq BA}
\]
\begin{block}{Message}
Matrix multiplication is associative and distributive, but not commutative in general.
\end{block}
\end{frame}

\begin{frame}{Matrix-Vector Products}
Let
\[
A=\begin{pmatrix}2&-1&3\\0&4&1\end{pmatrix},
\qquad
x=\begin{pmatrix}1\\-2\\0\end{pmatrix}.
\]
\pause
\[
Ax=
\begin{pmatrix}
2(1)+(-1)(-2)+3(0)\\
0(1)+4(-2)+1(0)
\end{pmatrix}
=
\begin{pmatrix}4\\-8\end{pmatrix}.
\]
\pause
\begin{block}{Interpretation}
Each component of $Ax$ is a dot product between a row of $A$ and the vector $x$.
\end{block}
\end{frame}

\section{Associated Operators: Transpose, Trace, and Determinant}

\begin{frame}{Transpose: Definition and Example}
For $A=(a_{ij})\in\M_{m,n}(\R)$, the transpose $A^\top\in\M_{n,m}(\R)$ is
defined by
\[
(A^\top)_{ij}=a_{ji}.
\]
It turns rows into columns and columns into rows.
\pause
\[
\begin{pmatrix}1&3&-2\\0&4&5\end{pmatrix}^{\!\top}
=\begin{pmatrix}1&0\\3&4\\-2&5\end{pmatrix}.
\]
\end{frame}

\begin{frame}{Transpose Rules and Dot Products}
Whenever dimensions are compatible,
\[
(A+B)^\top=A^\top+B^\top,
\qquad
(\lambda A)^\top=\lambda A^\top,
\qquad
(AB)^\top=B^\top A^\top.
\]
\pause
For vectors $x,y\in\R^d$,
\[
x^\top y=\sum_{i=1}^d x_i y_i,
\qquad
(Ax)^\top y=x^\top A^\top y.
\]
\begin{block}{Why the order reverses}
The product $AB$ means first apply $B$, then $A$. Transposition reverses this
composition, hence $(AB)^\top=B^\top A^\top$.
\end{block}
\end{frame}

\begin{frame}{Symmetric and Positive Semidefinite Matrices}
A is symmetric when $A^\top=A$. It is positive semidefinite when
\[
x^\top Ax\ge0\qquad\text{for every }x\in\R^n.
\]
\pause
For example,
\[
A=\begin{pmatrix}2&1\\1&2\end{pmatrix},
\qquad
x^\top Ax=(x_1+x_2)^2+x_1^2+x_2^2>0
\quad(x\neq0).
\]
\begin{block}{Why this matters}
Covariance matrices are symmetric positive semidefinite, and portfolio
variances are quadratic forms.
\end{block}
\end{frame}

\begin{frame}{Trace: Definition and Main Rules}
For $A=(a_{ij})\in\M_n(\R)$,
\[
\tr(A)=\sum_{i=1}^n a_{ii}.
\]
\pause
The trace is linear and cyclic:
\[
\tr(A+B)=\tr(A)+\tr(B),
\qquad
\tr(\lambda A)=\lambda\tr(A),
\qquad
\tr(AB)=\tr(BA).
\]
\pause
For example,
\[
\tr\begin{pmatrix}3&1&0\\2&-4&5\\1&1&7\end{pmatrix}=3-4+7=6.
\]
\end{frame}

\begin{frame}{Determinant: Meaning and the $2\times2$ Formula}
The determinant measures signed volume scaling and detects dimension loss.
For
\[
A=\begin{pmatrix}a&b\\c&d\end{pmatrix},
\qquad
\det(A)=ad-bc.
\]
\pause
\begin{block}{Invertibility criterion}
\[
A\text{ is invertible}\quad\Longleftrightarrow\quad\det(A)\neq0.
\]
If $\det(A)=0$, at least one direction is collapsed.
\end{block}
\pause
\[
\det\begin{pmatrix}2&3\\1&-4\end{pmatrix}=-11,
\qquad
\det\begin{pmatrix}1&2\\2&4\end{pmatrix}=0.
\]
\end{frame}

\begin{frame}{The Horrible General Formula for the Determinant}
For $A=(a_{ij})\in\M_n(\R)$,
\[
\boxed{\det(A)=\sum_{\sigma\in\mathfrak S_n}
\operatorname{sgn}(\sigma)\prod_{i=1}^n a_{i,\sigma(i)}}.
\]
\pause
Here $\mathfrak S_n$ contains all permutations of $\{1,\ldots,n\}$ and
\[
\operatorname{sgn}(\sigma)=(-1)^{N(\sigma)},
\quad
N(\sigma)=\#\{(i,j):i<j,\ \sigma(i)>\sigma(j)\}.
\]
\begin{alertblock}{Computational warning}
There are $n!$ signed products. The formula is fundamental for theory, but
elimination or factorization is preferable in practice.
\end{alertblock}
\end{frame}

\begin{frame}{Laplace Expansion Along a Row or a Column}
Delete row $i$ and column $j$ to obtain $A_{ij}$. Define
\[
M_{ij}=\det(A_{ij}),
\qquad
C_{ij}=(-1)^{i+j}M_{ij}.
\]
\pause
Expansion along a fixed row $i$ gives
\[
\boxed{\det(A)=\sum_{j=1}^n a_{ij}C_{ij}}
=\sum_{j=1}^n(-1)^{i+j}a_{ij}\det(A_{ij}).
\]
Expansion along a fixed column $j$ gives
\[
\boxed{\det(A)=\sum_{i=1}^n a_{ij}C_{ij}}
=\sum_{i=1}^n(-1)^{i+j}a_{ij}\det(A_{ij}).
\]
\begin{block}{Practical choice}
Expand along a row or column containing many zeros.
\end{block}
\end{frame}

\begin{frame}{Computing a $3\times3$ Determinant}
For
\[
A=\begin{pmatrix}a&b&c\\d&e&f\\g&h&i\end{pmatrix},
\]
expansion along the first row gives
\[
\det(A)=a(ei-fh)-b(di-fg)+c(dh-eg).
\]
\pause
For example,
\[
\det\begin{pmatrix}1&2&0\\-1&3&1\\2&0&4\end{pmatrix}
=12-2(-6)=24.
\]
The nonzero determinant proves that the matrix is invertible.
\end{frame}

\begin{frame}{Determinant Rules}
For square matrices of the same size,
\begin{itemize}
  \item $\det(AB)=\det(A)\det(B)$ and $\det(A^\top)=\det(A)$;
  \item $\det(A^{-1})=1/\det(A)$ when $A$ is invertible;
  \item swapping two rows changes the sign;
  \item multiplying one row by $\lambda$ multiplies the determinant by $\lambda$;
  \item adding a multiple of one row to another leaves it unchanged.
\end{itemize}
\pause
For a triangular matrix,
\[
\det(A)=\prod_{i=1}^n a_{ii}.
\]
\end{frame}

\section{Inversion and Linear Systems}

\begin{frame}{Identity and Inverse}
\begin{block}{Identity matrix}
\[
(I_n)_{ij}=
\begin{cases}
1,& i=j,\\
0,& i\neq j.
\end{cases}
\]
It satisfies $AI_n=A$ and $I_nA=A$ when dimensions are compatible.
\end{block}
\pause
\begin{block}{Inverse}
A square matrix $A\in \M_n(\R)$ is invertible if there exists $A^{-1}$ such that
\[
AA^{-1}=A^{-1}A=I_n.
\]
If $A$ is invertible, the system $Ax=b$ has the unique solution $x=A^{-1}b$.
\end{block}
\end{frame}

\begin{frame}{Inverse of a $2\times2$ Matrix}
For
\[
A=\begin{pmatrix}a&b\\c&d\end{pmatrix},
\]
the inverse exists exactly when $\det(A)=ad-bc\neq0$. Then
\[
\boxed{
A^{-1}=\frac{1}{\det(A)}
\begin{pmatrix}d&-b\\-c&a\end{pmatrix}
=\frac{1}{ad-bc}
\begin{pmatrix}d&-b\\-c&a\end{pmatrix}.}
\]
\pause
For the matrix used in the next example,
\[
\begin{pmatrix}2&3\\1&-4\end{pmatrix}^{-1}
=\begin{pmatrix}4/11&3/11\\1/11&-2/11\end{pmatrix}.
\]
\end{frame}

\begin{frame}{Worked Example: Solving a Linear System}
Consider the system $Ax=b$ with
\[
A=\begin{pmatrix}2&3\\1&-4\end{pmatrix},
\qquad b=\begin{pmatrix}5\\-2\end{pmatrix},
\qquad \det(A)=-11\neq0.
\]
\pause
Substitution gives
\[
11x_2=9,
\qquad x_2=\frac9{11},
\qquad x_1=-2+4x_2=\frac{14}{11}.
\]
\pause
\[
\boxed{x=\begin{pmatrix}14/11\\9/11\end{pmatrix}}
\]
Direct substitution and the inverse formula give the same solution.
\end{frame}

\begin{frame}{Triangular Systems and Forward Substitution}
Suppose
\[
L=\begin{pmatrix}
1&0&0\\
1&1&0\\
1&1&1
\end{pmatrix},
\qquad
Ly=
\begin{pmatrix}P_1\\P_2\\P_3\end{pmatrix}.
\]
\pause
The equations are
\[
y_1=P_1,\qquad y_1+y_2=P_2,\qquad y_1+y_2+y_3=P_3.
\]
\pause
\[
y_1=P_1,\qquad y_2=P_2-P_1,\qquad y_3=P_3-P_2.
\]
\end{frame}

\begin{frame}{Finance Example: Discount Factors}
Zero-coupon prices are observed:
\[
P_1=0.97,\qquad P_2=0.92,\qquad P_3=0.855.
\]
Write discount factors as partial sums:
\[
D(T_1)=y_1,\quad D(T_2)=y_1+y_2,\quad D(T_3)=y_1+y_2+y_3.
\]
\pause
\[
\begin{pmatrix}
1&0&0\\
1&1&0\\
1&1&1
\end{pmatrix}
\begin{pmatrix}y_1\\y_2\\y_3\end{pmatrix}
=
\begin{pmatrix}0.97\\0.92\\0.855\end{pmatrix}.
\]
\pause
\[
y_1=0.97,\qquad y_2=-0.05,\qquad y_3=-0.065.
\]
\end{frame}

\section{Injectivity, Surjectivity, Image, Kernel, and Rank}

\begin{frame}{Image and Kernel}
\small
Let $A\in\M_{m,n}(\R)$, viewed as a linear map from $\R^n$ to $\R^m$.
\begin{block}{Image / range}
\[
\im(A)=\{Ax:x\in\R^n\}\subset \R^m.
\]
It is the subspace reached by $A$.
\end{block}
\pause
\begin{block}{Kernel / null space}
\[
\ker(A)=\{x\in\R^n:Ax=0\}.
\]
It is the set of directions collapsed by $A$.
\end{block}
\pause
\begin{resultblock}{Result -- Rank--nullity theorem}
The rank is $\rank(A)=\dim\im(A)$, and
\[
\rank(A)+\dim\ker(A)=n.
\]
\end{resultblock}
\end{frame}

\begin{frame}{Invertibility Criteria}
\begin{resultblock}{Result -- Equivalent invertibility criteria}
For $A\in \M_n(\R)$, the following statements are equivalent:
\begin{itemize}
  \item $A$ is invertible;
  \item $\ker(A)=\{0\}$;
  \item $\im(A)=\R^n$;
  \item $\rank(A)=n$;
  \item $Ax=b$ has a unique solution for every $b\in\R^n$.
\end{itemize}
\end{resultblock}
\pause
\begin{alertblock}{Practical warning}
For numerical work, we usually solve $Ax=b$ directly instead of computing $A^{-1}$ explicitly.
\end{alertblock}
\end{frame}

\begin{frame}{Orthogonality and Orthogonal Complements}
Two vectors $x,y\in\R^m$ are orthogonal when
\[
x\perp y
\quad\Longleftrightarrow\quad
x^\top y=0.
\]
\pause
For a subspace $V\subset\R^m$, define
\[
\boxed{V^\perp
=\{y\in\R^m:y^\top v=0\text{ for every }v\in V\}.}
\]
\begin{block}{Geometric meaning}
$V^\perp$ contains all directions perpendicular to every vector of $V$.
\end{block}
\pause
For $V=\operatorname{span}\{(1,1,1)^\top\}$,
\[
V^\perp=\{y\in\R^3:y_1+y_2+y_3=0\}.
\]
\end{frame}

\begin{frame}{Image, Kernel, and Orthogonality}
For $A\in \M_{m,n}(\R)$,
\[
\im(A^\top)=(\ker A)^\perp,
\qquad
\im(A)=(\ker A^\top)^\perp.
\]
\pause
\begin{block}{Interpretation}
The row space contains exactly the directions that are visible to the linear map.
The kernel contains directions that disappear.
\end{block}
\end{frame}

\begin{frame}{Example: Kernel and Image}
Let
\[
A=\begin{pmatrix}
1&-1&0\\
-1&2&-1\\
0&-1&1
\end{pmatrix}.
\]
\pause
Solving $Ax=0$ gives
\[
x_1=x_2=x_3,\qquad
\ker(A)=\operatorname{span}\left\{\begin{pmatrix}1\\1\\1\end{pmatrix}\right\}.
\]
\pause
Since $\dim\ker(A)=1$ and $A$ has three columns,
\[
\rank(A)=2.
\]
Because $A$ is symmetric, $\im(A)=(\ker A)^\perp$ is the plane
$y_1+y_2+y_3=0$.
\end{frame}

\section{Classical Decompositions}

\subsection{Eigenvalue Diagonalization}

\begin{frame}{Eigenvalues and Eigenvectors}
An eigenvector of $A\in \M_n(\R)$ is a nonzero vector $v$ such that
\[
Av=\lambda v.
\]
The scalar $\lambda$ is the associated eigenvalue.
\pause
\begin{block}{Geometric meaning}
In an eigenvector direction, the matrix acts by stretching, shrinking, or reversing direction. The direction is preserved.
\end{block}
\pause
Eigenvalues are found by solving
\[
\det(A-\lambda I)=0.
\]
\end{frame}

\begin{frame}{Diagonalization}
A matrix $A\in\M_n(\R)$ is diagonalizable if there exist an invertible matrix $P$ and a diagonal matrix $D$ such that
\[
A=PDP^{-1}.
\]
\pause
The columns of $P$ are eigenvectors and the diagonal entries of $D$ are eigenvalues:
\[
P=\begin{pmatrix}|&&|\\v_1&\cdots&v_n\\|&&|\end{pmatrix},
\qquad
D=\diag(\lambda_1,\ldots,\lambda_n).
\]
\pause
\begin{resultblock}{Result -- Powers of a diagonalizable matrix}
\[
A^k=PD^kP^{-1},\qquad
D^k=\diag(\lambda_1^k,\ldots,\lambda_n^k).
\]
\end{resultblock}
\end{frame}

\begin{frame}{Diagonalization: Algorithm}
\begin{enumerate}
  \item Compute the characteristic polynomial $p_A(\lambda)=\det(A-\lambda I)$.
  \item Find the eigenvalues.
  \item For each eigenvalue, solve $(A-\lambda I)v=0$.
  \item If there are $n$ linearly independent eigenvectors, build $P$.
  \item Write $A=PDP^{-1}$.
\end{enumerate}
\pause
\begin{alertblock}{Warning}
Not every matrix is diagonalizable. A repeated eigenvalue may fail to provide enough independent eigenvectors.
\end{alertblock}
\end{frame}

\begin{frame}{Diagonalization: Worked Example}
Let
\[
A=\begin{pmatrix}4&1\\2&3\end{pmatrix}.
\]
\pause
The characteristic polynomial is
\[
\det(A-\lambda I)=(\lambda-5)(\lambda-2),
\qquad \lambda_1=5,\quad \lambda_2=2.
\]
\pause
Eigenvectors:
\[
\lambda_1=5:\quad v_1=\begin{pmatrix}1\\1\end{pmatrix},
\qquad
\lambda_2=2:\quad v_2=\begin{pmatrix}1\\-2\end{pmatrix}.
\]
\pause
\[
P=\begin{pmatrix}1&1\\1&-2\end{pmatrix},
\qquad D=\begin{pmatrix}5&0\\0&2\end{pmatrix},
\qquad
A=PDP^{-1}.
\]
\end{frame}

\begin{frame}{Triangular Example for the Dynamics}
For the dynamical application, consider
\[
A=\begin{pmatrix}3&1\\0&2\end{pmatrix}.
\]
Its eigenvalues are the diagonal entries, with eigenvectors
\[
v_1=\begin{pmatrix}1\\0\end{pmatrix},
\qquad
v_2=\begin{pmatrix}-1\\1\end{pmatrix}.
\]
Hence
\[
P=\begin{pmatrix}1&-1\\0&1\end{pmatrix},
\qquad D=\begin{pmatrix}3&0\\0&2\end{pmatrix},
\qquad A=PDP^{-1}.
\]
\end{frame}

\begin{frame}{Application: A Discrete-Time Dynamical System}
Consider $x_{k+1}=Ax_k$ with
\[
A=\begin{pmatrix}3&1\\0&2\end{pmatrix},
\qquad x_0=\begin{pmatrix}0\\1\end{pmatrix}.
\]
Using the previous diagonalization,
\[
A^k=PD^kP^{-1}
=\begin{pmatrix}
3^k&3^k-2^k\\
0&2^k
\end{pmatrix}.
\]
\pause
Therefore
\[
x_k=A^kx_0
=\begin{pmatrix}3^k-2^k\\2^k\end{pmatrix}.
\]
\begin{block}{Interpretation}
Diagonalization separates the dynamics into two scalar growth modes, $3^k$ and $2^k$.
\end{block}
\end{frame}

\subsection{Spectral Decomposition}

\begin{frame}{Symmetric Matrices: Spectral Theorem}
\begin{resultblock}{Result -- Spectral theorem}
If $A$ is real and symmetric, then:
\begin{itemize}
  \item all eigenvalues are real;
  \item eigenvectors associated with distinct eigenvalues are orthogonal;
  \item there exists an orthogonal matrix $Q$ and a diagonal matrix $D$ such that
  \[
  A=QDQ^\top.
  \]
\end{itemize}
\end{resultblock}
\pause
\begin{block}{Why this is central}
Covariance matrices are symmetric. Their eigenvectors define orthogonal risk factors or principal components.
\end{block}
\end{frame}

\begin{frame}{Application: Common and Spread Risk Factors}
Consider the covariance matrix
\[
C=\begin{pmatrix}1&0.8\\0.8&1\end{pmatrix}.
\]
Its spectral decomposition has
\[
\lambda_1=1.8,\quad q_1=\frac1{\sqrt2}\begin{pmatrix}1\\1\end{pmatrix},
\qquad
\lambda_2=0.2,\quad q_2=\frac1{\sqrt2}\begin{pmatrix}1\\-1\end{pmatrix}.
\]
\pause
The common position $q_1$ has variance $1.8$, whereas the long--short position
$q_2$ has variance $0.2$.
\begin{block}{Interpretation}
The spectral theorem identifies orthogonal portfolios and quantifies the risk carried by each factor.
\end{block}
\end{frame}

\subsection{Cholesky Decomposition}

\begin{frame}{Cholesky Decomposition}
\begin{resultblock}{Result -- Cholesky decomposition}
If $A$ is symmetric positive definite, there is a unique lower triangular
matrix $L$ with positive diagonal entries such that
\[
\boxed{A=LL^\top.}
\]
\end{resultblock}
\pause
For the positive-definite matrix
\[
A=\begin{pmatrix}4&2\\2&3\end{pmatrix},
\qquad
L=\begin{pmatrix}2&0\\1&\sqrt2\end{pmatrix},
\qquad
A=LL^\top.
\]
\pause
\begin{block}{Why this matters}
Cholesky is used to solve positive-definite systems and to simulate correlated
Gaussian vectors: if $Z\sim N(0,I)$ and $\Sigma=LL^\top$, then $LZ$ has
covariance $\Sigma$.
\end{block}
\end{frame}

\begin{frame}{Application: Simulating Correlated Returns}
Let $Z_1,Z_2$ be independent standard normal variables and set
\[
L=\begin{pmatrix}2&0\\1&\sqrt2\end{pmatrix},
\qquad
X=L\begin{pmatrix}Z_1\\Z_2\end{pmatrix}
=\begin{pmatrix}2Z_1\\Z_1+\sqrt2 Z_2\end{pmatrix}.
\]
\pause
Then
\[
\Cov(X)=LL^\top
=\begin{pmatrix}4&2\\2&3\end{pmatrix}.
\]
Thus $X_1$ and $X_2$ have variances $4$ and $3$, covariance $2$, and
correlation $2/\sqrt{12}=1/\sqrt3$.
\begin{block}{Practical use}
Generate independent shocks first, then multiply by $L$ to obtain shocks with the desired covariance.
\end{block}
\end{frame}

\section{Statistics and Finance Applications}

\begin{frame}{Covariance Matrix from Data}
Let $X\in\R^{n\times d}$ be a centered data matrix:
\[
X=
\begin{pmatrix}
- & x_1^\top & -\\
- & x_2^\top & -\\
&\vdots&\\
- & x_n^\top & -
\end{pmatrix}.
\]
\pause
The empirical covariance matrix is
\[
\widehat{\Sigma}=\frac{1}{n-1}X^\top X.
\]
\pause
\begin{resultblock}{Result -- Covariance matrices are positive semidefinite}
\[
\widehat{\Sigma}^\top=\widehat{\Sigma},
\qquad
z^\top\widehat{\Sigma}z=\frac{1}{n-1}\|Xz\|^2\ge 0.
\]
\end{resultblock}
\end{frame}

\begin{frame}{Covariance Matrix: Numerical Example}
Use the centered data matrix
\[
X=\begin{pmatrix}
0.01&-0.03\\
0.02&0.01\\
-0.01&0
\end{pmatrix}.
\]
\pause
Then
\[
X^\top X=
\begin{pmatrix}
0.0006&-0.0001\\
-0.0001&0.0010
\end{pmatrix},
\qquad
\widehat\Sigma=\frac1{n-1}X^\top X.
\]
The matrix is symmetric and positive semidefinite by construction.
\end{frame}

\begin{frame}{Portfolio Variance}
\begin{resultblock}{Result -- Portfolio variance formula}
Let $w$ be a vector of portfolio weights and $\Sigma$ a covariance matrix of asset returns.
\[
R_p=w^\top r,\qquad
\boxed{\Var(R_p)=w^\top\Sigma w.}
\]
\end{resultblock}
\pause
Example:
\[
\Sigma=
\begin{pmatrix}
0.04&0.018\\
0.018&0.09
\end{pmatrix},
\qquad
w=\begin{pmatrix}0.6\\0.4\end{pmatrix}.
\]
\pause
\[
w^\top\Sigma w
=0.6^2(0.04)+2(0.6)(0.4)(0.018)+0.4^2(0.09)
=0.03744.
\]
\end{frame}

\begin{frame}{PCA in One Slide}
For centered data $X$, PCA diagonalizes
\[
\widehat{\Sigma}=\frac{1}{n-1}X^\top X.
\]
\pause
If
\[
\widehat{\Sigma}=QDQ^\top,
\]
then:
\begin{itemize}
  \item columns of $Q$ are principal directions;
  \item diagonal entries of $D$ are variances along those directions;
  \item the first principal component is the direction of maximum variance.
\end{itemize}
\pause
Projection onto the first $k$ components:
\[
Z=XQ_k.
\]
\end{frame}

\begin{frame}{PCA Example}
Centered data:
\[
X=
\begin{pmatrix}
-1&-1\\
0&0\\
1&1
\end{pmatrix}.
\]
\pause
\[
\widehat{\Sigma}
=\frac{1}{2}X^\top X
=
\begin{pmatrix}1&1\\1&1\end{pmatrix}.
\]
\pause
Eigenvalues and principal directions:
\[
\lambda_1=2,\quad q_1=\frac{1}{\sqrt 2}\begin{pmatrix}1\\1\end{pmatrix},
\qquad
\lambda_2=0,\quad q_2=\frac{1}{\sqrt 2}\begin{pmatrix}1\\-1\end{pmatrix}.
\]
All variance lies on the line $x_1=x_2$.
\end{frame}

\section[Markov Chains]{Markov Chains and Transition Matrices}

\subsection{Transition Matrices}

\begin{frame}{From a Markov Chain to a Matrix}
For states $\{1,\ldots,m\}$, use the column-vector convention
\[
P_{ij}=\mathbb{P}(X_{n+1}=i\mid X_n=j).
\]
\pause
Each column describes the next-state distribution from one current state:
\[
P_{ij}\geq0,
\qquad
\sum_{i=1}^mP_{ij}=1.
\]
\begin{block}{Linear algebra link}
A transition matrix is a nonnegative matrix that maps probability vectors to
probability vectors.
\end{block}
\end{frame}

\subsection{Distribution Dynamics}

\begin{frame}{Matrix Powers Represent Time}
\begin{resultblock}{Result -- Distribution dynamics}
Let $\mu_n$ be the column vector containing the distribution of $X_n$. Then
\[
\boxed{\mu_{n+1}=P\mu_n},
\qquad
\boxed{\mu_n=P^n\mu_0}.
\]
\end{resultblock}
\pause
The entry $(P^n)_{ij}$ is the probability of reaching state $i$ after $n$
steps when starting from state $j$.
\begin{block}{Key message}
Matrix multiplication composes transitions; taking the $n$th power advances
the chain by $n$ time steps.
\end{block}
\end{frame}

\begin{frame}{Example: A Two-State Credit Chain}
States are Good (G) and Bad (B), with
\[
P=\begin{pmatrix}0.9&0.4\\0.1&0.6\end{pmatrix},
\qquad
\mu_0=\begin{pmatrix}1\\0\end{pmatrix}.
\]
\pause
After one and two years,
\[
\mu_1=P\mu_0=\begin{pmatrix}0.9\\0.1\end{pmatrix},
\qquad
\mu_2=P^2\mu_0=\begin{pmatrix}0.85\\0.15\end{pmatrix}.
\]
\begin{block}{Interpretation}
Starting from Good, the two-year probability of being in Bad is $15\%$.
\end{block}
\end{frame}

\subsection{Stationarity and Diagonalization}

\begin{frame}{Stationary Distribution as an Eigenvector}
\small
\begin{resultblock}{Result -- Stationarity condition}
A stationary distribution satisfies
\[
P\pi=\pi,
\qquad
\mathbf{1}^\top\pi=1.
\]
Thus $\pi$ is an eigenvector for eigenvalue $1$. For the credit chain,
\end{resultblock}
\[
\pi=\begin{pmatrix}0.8\\0.2\end{pmatrix}.
\]
\pause
The other eigenvalue is $0.5$, and for $\mu_0=(1,0)^\top$,
\[
\mu_n=
\begin{pmatrix}0.8\\0.2\end{pmatrix}
+0.2(0.5)^n\begin{pmatrix}1\\-1\end{pmatrix}.
\]
The stationary mode persists, while the transient mode decays geometrically.
\end{frame}

\end{document}
